% =============================================================================
% ConfPert Phase 2 paper -- main.tex
% Target: ICLR 2027 Datasets & Benchmarks (Sept-Oct 2026 deadline)
% Template: NeurIPS D&B 2025 placeholder; swap to ICLR 2027 sty when CFP opens.
% Compile: tectonic -X compile main.tex --reruns 4
%
% Rewrite rules applied (per the strict-prompt protocol):
%   - [CITATION NEEDED] for unverified claims
%   - [NEEDS USER INPUT] for missing experimental detail
%   - Soften "first/novel/state-of-the-art" unless evidence supports
%   - Methods reproducibility: dataset version + seed + hardware + commit SHA
%   - Cite Csendes 2025 only for the empirical finding it actually reports
%     (FMs lose to mean baseline on Perturb-seq), not for an "embedding-baseline
%     mode" methodological recipe.
%   - the abstract opening verb is used only there (no repetition elsewhere).
% =============================================================================

\documentclass[11pt]{article}

\usepackage[a4paper, margin=1in]{geometry}
\usepackage{natbib}
\usepackage[utf8]{inputenc}
\usepackage{booktabs}
\usepackage{amsmath}
\usepackage{amssymb}
\usepackage{graphicx}
\usepackage{hyperref}
\usepackage{xcolor}
\usepackage{nicefrac}
\usepackage{microtype}

\newcommand{\TODO}[1]{{\color{red}\textbf{TODO:} #1}}
\newcommand{\citeneeded}{{\color{orange}\textsf{[CITATION NEEDED]}}}
\newcommand{\userinput}{{\color{orange}\textsf{[NEEDS USER INPUT]}}}

\title{Capacity is not Calibration:\\
A Pre-Registered Distribution-Aware Benchmark of Single-Cell Perturbation Predictors}

\author{
  Aayan Alwani \\
  Independent Researcher \\
  \texttt{alwaniaayan6@gmail.com}
}

\begin{document}
\maketitle

% =============================================================================
% Abstract  (~200 words; written last, per strict-prompt STEP 6)
% =============================================================================
\begin{abstract}
We introduce ConfPert, a benchmark that evaluates 12 single-cell
perturbation predictors with finite-sample conformal coverage as the
primary outcome rather than mean-prediction accuracy. The benchmark
covers 8 publicly available Perturb-seq and chemical-perturbation
datasets and ships with a machine-verifiable YAML pre-registration that
a command-line verifier refuses to evaluate when the locked
hypotheses, thresholds, or BH-FDR corrections do not match the results
file. Across 651 locked $(\text{predictor}, \text{dataset}, \alpha,
\text{discrepancy})$ cells, the pre-registered family-level
Kruskal-Wallis test ($H_3$) rejects equality of calibration deviation
between the baseline family ($F_1$, mean / bilinear-ridge / noisy-mean;
$\bar d = 0.078$, $n = 25$) and the transformer-foundation-model family
($F_3$, scGPT / scFoundation / Geneformer in frozen-encoder embedding
mode; $\bar d = 0.154$, $n = 25$) at
$p_\text{KW} = 1.7 \times 10^{-7}$ with Cliff's $\delta = 0.80$. The
hypotheses on per-dataset capacity ($H_1$) and cell-line context ($H_2$)
also pass; the hypotheses on training-corpus size ($H_{1b}$) and
corpus-diversity ($H_{2b}$) fail. The result reproduces the perturbation
prediction skepticism of \cite{ahlmann2025deep,csendes2025train} under a
coverage criterion. We release the library on PyPI
(\texttt{confpert v0.2.0rc2}), the pre-registration template, and a
Codabench leaderboard for community submissions.
\end{abstract}

% =============================================================================
% 1. Introduction  (~3/4 page)
% =============================================================================
\section{Introduction}

Predicting single-cell responses to genetic or chemical perturbations is
the headline application of single-cell foundation models. Recent
benchmarks report that simple linear baselines match or exceed deep
models on standard transcriptomic prediction metrics
\cite{ahlmann2025deep,csendes2025train,kedzierska2025zeroshot}, but the
metrics in these benchmarks are still primarily measures of mean-vector
accuracy (MSE, Pearson, top-$k$ DEG recovery). For downstream uses such
as cell-population sampling or risk-aware drug response, coverage of the
predicted distribution -- not mean accuracy -- is the actionable
property.

We evaluate predictors by per-population conformal coverage. Six
discrepancy scores -- Kolmogorov-Smirnov, 1-Wasserstein, energy,
MMD-RBF, bimodality mismatch, and variance-ratio deviation -- are
wrapped in finite-sample conformal heads \cite{vovk2012,tibshirani2019,
romano2019,izbicki2022,cauchois2021,bates2021}, and we report
calibration deviation
$|{\hat{\alpha}} - \alpha|$ as the primary outcome. The pipeline is
locked in a YAML pre-registration before any first run; a CLI verifier
(\texttt{confpert prereg verify}) refuses to evaluate when the YAML and
the results file disagree. Pre-registration in machine learning has
been argued for \cite{hofman2023prereg,pineau2020reproducibility} but
remains rare; the closest mechanically enforced precedents we are aware
of are competition-style leaderboards (e.g.\ Kaggle), not
hypothesis-level locks. \citeneeded{} for a single-cell-perturbation
benchmark with a comparable mechanically enforced lock.

\noindent\emph{Empirical finding.}
The pre-registered Kruskal-Wallis test ($H_3$) rejects equality of
calibration deviation across the three predictor families
($p_\text{KW} = 1.7 \times 10^{-7}$, Cliff's $\delta = 0.80$). The
baseline family ($F_1$) is the best-calibrated; the
foundation-model family ($F_3$) is the worst. Per-dataset capacity
($H_1$, 4 of 8 datasets at $\rho > 0.5,\, p < 0.05$) and cell-line
context ($H_2$, two-way ANOVA $f$-$p = 0.0064$, $\eta^2 = 0.120$) also
pass; corpus-size and corpus-diversity hypotheses ($H_{1b}$, $H_{2b}$)
fail.

\noindent\emph{Contributions.}
\begin{itemize}
  \item A benchmark of 12 published predictors across 8 public
        single-cell perturbation datasets with within-perturbation and
        cross-cell-line splits; 651 locked result cells.
  \item A machine-verifiable YAML pre-registration with a CLI verifier
        that refuses to run when the locked YAML and the results file
        disagree.
  \item An ablation panel that re-runs the same hypothesis tests under
        counterfactual analyst-side relaxations, quantifying the
        threshold-shopping bias avoided by the lock.
  \item A pip-installable library, datasheet and model card, and a
        Codabench leaderboard.
\end{itemize}

% =============================================================================
% 2. Related work  (~1/2 page)
% =============================================================================
\section{Related work}

\paragraph{Conformal prediction for multivariate and generative outputs.}
Conformal prediction has been extended to the multivariate and
generative-model setting in a recent cluster of methods
\cite{meyer2026kernel,zheng2024generative,thurin2025optimal,
ndiaye2025multivariate,dheur2025generalized,yang2026cp4gen,
braun2025minvol,otcp2025}, including a perturbation-data application
\cite{asiaee2026partial}. ConfPert sits downstream of this literature
and operationalises six standard per-population discrepancies inside
finite-sample conformal heads.

\paragraph{Single-cell foundation models and perturbation prediction
benchmarks.}
Several recent benchmarks evaluate single-cell foundation models on
perturbation tasks. \citet{ahlmann2025deep} compare five foundation
models and two other deep learning models against deliberately simple
baselines on Perturb-seq and report that none outperformed the
baselines. \citet{csendes2025train} replicate the finding for scGPT and
scFoundation against a mean-of-training-examples baseline.
\citet{kedzierska2025zeroshot} show that scGPT and Geneformer can
underperform simpler methods in zero-shot deployments.
\citet{boiarsky2024deeper} reach a similar conclusion under a deeper
single-cell foundation-model evaluation. \citet{li2025systematic}
surveys the largest pool of perturbation predictors to date but reports
accuracy-style scores (MSE, mean Pearson) rather than calibration
coverage. ConfPert's incremental contribution is to add a coverage
criterion, a finite-sample conformal head per discrepancy, and a
mechanically enforced pre-registration.

\paragraph{Pre-registration in machine learning.}
Pre-registration is established in social and biomedical research
\cite{nosek2018prereg} and has been argued for in predictive modelling
\cite{hofman2023prereg,pineau2020reproducibility}. The
\texttt{ConfPert} verifier extends this prior work by refusing to
evaluate when the locked YAML and the on-disk results file disagree.

% =============================================================================
% 3. Benchmark design  (~1.5 pages, the contribution chapter)
% =============================================================================
\section{Benchmark design}
\label{sec:design}

\subsection{Predictors}
\label{sec:predictors}
We benchmark 12 predictors across three pre-registered family levels.
$F_1$ baselines: mean (single Perturb-seq mean per perturbation),
ahlmann\_bilinear\_ridge \cite{ahlmann2025deep} (additive model with
ridge regularisation, $K = 10$, $\lambda = 0.1$), and noisy\_mean
(mean with isotropic noise). $F_2$ latent-delta: scGen
\cite{lotfollahi2019scgen}, biolord \cite{piran2024biolord}, CPA
\cite{lotfollahi2023cpa}, sVAE\textsuperscript{+}
\cite{bereket2023}, and STATE \cite{adduri2025state}. $F_3$
transformer: GEARS \cite{roohani2024gears}, scGPT
\cite{cui2024scgpt}, scFoundation \cite{hao2023scfoundation}, and
Geneformer \cite{theodoris2023geneformer}.

Foundation models run in frozen-encoder mode. A ridge regression head
(\texttt{sklearn.linear\_model.Ridge}, $\alpha_\text{ridge} = 1.0$) is
fitted on the train fold of cell-level encoder embeddings against the
log1p HVG expression target; the per-gene sampling variance used in the
predictive distribution is computed only on training-fold rows via
\texttt{confpert.heavyweight\_helpers.safe\_train\_variance}, which
raises if the supplied train and test row indices overlap. This
test-fold isolation is enforced as code so that a future contributor
cannot silently re-introduce the leak. Full fine-tuning of the
foundation models is left to future work; the present results should be
read as upper bounds on the F3 family under transfer-learning use, not
as a verdict on the architectures.

\subsection{Datasets}
\label{sec:datasets}
We use 8 datasets covered by at least one filled result row:
Norman~2019 \cite{norman2019}, Replogle K562 and RPE1 essential
\cite{replogle2022}, Adamson~2016 \cite{adamson2016}, Tahoe-100M
subset \cite{tahoe2025}, Frangieh \cite{frangieh2021melanoma},
Schmidt \cite{schmidt2022crispr}, Datlinger \cite{datlinger2017crop},
and Lara-Astiaso~2023 ex-vivo mouse Perturb-seq \userinput{} for full
in-text citation of the Lara-Astiaso 2023 Nature Genetics primary
source. McFaline-Figueroa \cite{mcfaline2024sciplex} and CellFM
\cite{zeng2024cellfm} are documented in the pre-registration but report
as N/A in the result table (sample-meta hash-table demultiplexing for
McFaline and absence of a usable HF mirror for CellFM as of
2026-05-25). The mouse-organism slot was originally pre-registered for
``Walker et al.\ Nature Methods 2022 BMDC Perturb-seq''; that
citation could not be cross-verified, and the slot was substituted with
Lara-Astiaso 2023 on 2026-05-25 with an explicit amendment block in
the YAML lock (option B of the three documented resolution paths in
\texttt{confpert/data/walker.py}).

\subsection{Splits}
Within-perturbation: each perturbation's cells are partitioned 50/25/25
into train/calibration/test by seed 42; control cells are split
proportionally. Held-out-perturbation: 50/50 perturbation-level split
by the same seed. Cross-cell-line: train on Replogle K562, evaluate on
Replogle RPE1 with gene-name intersection. Held-out-organism (train
Lara-Astiaso, test Norman), held-out-modality (Frangieh RNA train,
Frangieh protein test), and held-out-perturbation-type (Replogle K562
train, Tahoe chemical test) splits are pre-registered but the
held-out-modality and held-out-perturbation-type cells report as
\userinput{} in the present run (those split-specific pipelines were
not executed within the Modal budget).

\subsection{Discrepancies and conformal heads}
Six per-population discrepancies: Kolmogorov-Smirnov on per-gene
marginals, 1-Wasserstein on per-gene marginals, energy distance
\cite{szekely2013}, MMD with an RBF kernel \cite{romano2019,
izbicki2022}, bimodality mismatch (KDE-based per-gene bimodality
divergence), and variance-ratio deviation
(\texttt{var(pred) / var(obs) - 1}). Each discrepancy is wrapped in a
per-perturbation conformal head
\cite{vovk2012,tibshirani2019,romano2019}; subgroup-conditional
\cite{cauchois2021} and RCPS \cite{bates2021} heads are implemented in
the library but evaluated as \userinput{} in this paper.

\subsection{Pre-registration protocol}
The hypotheses, thresholds, BH-FDR corrections, and disposition rules
are committed to \texttt{paper\_neurips\_dnb\_2026/preregistration\_v2.yaml}
before any first run; the file's git commit SHA and SHA-256 hash are
written into the YAML's \texttt{lock} block. \texttt{confpert prereg
verify} refuses to run if the on-disk YAML hash differs from the locked
hash, if the git commit SHA on the YAML differs from the locked SHA, or
if the \texttt{results\_json\_max\_mtime} field is exceeded by the
results file's mtime. Re-emitting hashes
(\texttt{confpert prereg emit-hashes}) is required after any
intentional amendment and is itself recorded by a git commit.

% =============================================================================
% 4. Materials and Methods  (strict-prompt STEP 2: reproducibility-grade)
% =============================================================================
\section{Materials and Methods}
\label{sec:methods}

\noindent\emph{Software environment.}
Python 3.11 (CPU) and 3.10 (CPA / scvi-tools subset). Key version pins:
\texttt{torch 2.4.1} (\texttt{torch 2.3.1} in the scGPT image to keep
\texttt{torchtext 0.18.0} ABI-compatible), \texttt{scvi-tools<1.4} for
scGen / biolord, \texttt{cell-gears 0.1.2} for GEARS,
\texttt{scgpt 0.2.4}, \texttt{transformers>=4.40,<4.46}. The full image
graph (one per heavyweight predictor family) is defined in
\texttt{scripts/modal\_launch.py}. \userinput{} for the exact Docker
image SHA when frozen for camera-ready.

\noindent\emph{Hardware.}
Heavyweight predictor runs were executed on Modal A100 GPUs (one A100
per (predictor, dataset) call, 24--32~GB host RAM); lightweight predictor
runs on a local M-series Apple Silicon CPU. The Modal budget cap was
\$3000; the actual total spend on the locked sweep was \userinput{} for
the camera-ready figure.

\noindent\emph{Datasets and preprocessing.}
The eight datasets and their source URLs are listed in the YAML
(\texttt{k1\_v2.datasets}) and reproduced in
Appendix~\ref{sec:datasheet}. All datasets are loaded by the
common pipeline (\texttt{confpert.data.\_load\_perturb\_ds}): drop
cells with fewer than 200 UMI; drop genes detected in fewer than 5
cells; \texttt{scanpy.pp.normalize\_total(target\_sum=1e4)};
\texttt{scanpy.pp.log1p}; \texttt{scanpy.pp.highly\_variable\_genes}
with $n_\text{top} = 512$. The top 50 perturbations by cell count are
retained (29 for Datlinger). Random seed is 42 throughout.

\noindent\emph{Predictor fitting.}
$F_1$ predictors are fit on the train fold and noise variants A--D
sweep isotropic, per-gene marginal, and full-covariance Gaussian noise.
$F_2$ and $F_3$ predictors are fit per the upstream package defaults;
all per-predictor configurations are stored in
\texttt{baselines/results.json:rows[*].config}. Foundation models pass
through \texttt{safe\_train\_variance} and
\texttt{sample\_from\_train\_distribution}
(\texttt{src/confpert/heavyweight\_helpers.py}); both refuse mis-split
inputs.

\noindent\emph{Discrepancy and conformal-head computation.}
For each held-out perturbation $p$, we draw a predicted population
$\hat{X}_p$ of size matching the observed cell count, compute each of
the six discrepancies $s_i(\hat{X}_p, X_p)$, and run a finite-sample
per-perturbation conformal calibration on the train half of the
perturbations. The resulting threshold $\tau_\alpha$ is applied to the
held-out perturbations and we report achieved coverage
$\hat\alpha = \Pr[s_i \leq \tau_\alpha]$.

\noindent\emph{Hypothesis tests.}
$H_1$: Spearman $\rho$ between $\log_{10}(\text{param count})$ and mean
calibration deviation per dataset; pass if $\rho > 0.5$, $p < 0.05$
(permutation, $n = 10000$), in $\geq 3$ of 4 main datasets.
$H_{1b}$: Spearman with training-corpus size, $\rho < -0.5$, same
$\geq 3$ of 4 rule.
$H_2$: Two-way Type-II ANOVA on calibration deviation with factors
(cell-line context, $\log$ param count); pass if main-effect $f$-$p <
0.0167$ (Bonferroni on 3 family tests) and $\eta^2 > 0.10$.
$H_{2b}$: Spearman with $\log$ corpus diversity, $\rho < -0.5$,
$p < 0.0167$, $\geq 7$ of 10 datasets.
$H_3$: Kruskal-Wallis across the three families, $p < 0.0167$ and
Cliff's $\delta > 0.50$ between best and worst.

\noindent\emph{Bootstrap CIs and power.}
Wilson 95\% binomial intervals are computed per
(predictor, dataset, score, $\alpha$) cell; per-cell counts and
intervals are reported in
\texttt{phase2d\_report\_phase2.json:bootstrap\_ci\_table}. Power
estimates use Monte Carlo with $n_\text{sims} = 2000$ for the
Spearman / Kruskal tests and \texttt{statsmodels.stats.power.FTestAnovaPower}
for the ANOVA.

\noindent\emph{Reproduction.}
\begin{verbatim}
git checkout <commit-SHA>
pip install -e .
python -m confpert.cli phase2d \
  --prereg paper_neurips_dnb_2026/preregistration_v2.yaml \
  --results baselines/results.json \
  --out paper_neurips_dnb_2026/phase2d_report_phase2.json
\end{verbatim}

% =============================================================================
% 5. Results (strict-prompt STEP 3: numbers from verified outputs)
% =============================================================================
\section{Results}
\label{sec:results}

\subsection{K1 sweep}
\label{sec:k1}
The K1 sweep produced 651 cells across 12 predictors, 8 datasets, 3
target coverage levels ($\alpha \in \{0.05, 0.10, 0.20\}$), and 6
discrepancies. Mean calibration deviation per family is
$\bar d_{F_1} = 0.078$ ($n = 25$ cells), $\bar d_{F_2} = 0.123$
($n = 26$), and $\bar d_{F_3} = 0.154$ ($n = 25$). The
per-discrepancy hardness rank (mean deviation across all
predictor-dataset pairs, ascending) is KS $<$ bimodality\_mismatch $<$
MMD-RBF $<$ $W_1$ $<$ energy $<$ variance\_ratio\_dev. The cross-dataset
heatmap is in Figure~\ref{fig:k1heatmap} and the
calibration-vs-capacity Pareto is in Figure~\ref{fig:pareto}; full
per-cell results with Wilson 95\%~CIs are in Appendix~\ref{sec:k1full}
and the JSON sidecar.

\subsection{K2 v2 hypothesis tests}
\label{sec:k2v2}
The five pre-registered hypotheses resolved as follows
(Table~\ref{tab:disp}).
$H_1$ \emph{capacity} passes: 4 of 8 datasets (Norman, Frangieh,
Replogle~RPE1, Tahoe) meet the per-dataset $\rho > 0.5$, $p < 0.05$
rule. $H_{1b}$ \emph{data} fails: per-dataset Spearman is degenerate
because many predictors share corpus size. $H_2$ \emph{cell-line
context} passes: two-way ANOVA $f$-$p = 0.0064$ ($< 0.0167$) and
$\eta^2 = 0.120$ ($> 0.10$). $H_{2b}$ \emph{corpus diversity} fails:
$\rho > 0$ on every dataset (opposite sign of the hypothesised
$\rho < -0.5$). $H_3$ \emph{architecture family} passes: Kruskal-Wallis
$p = 1.7 \times 10^{-7}$ with Cliff's $\delta = 0.798$ between $F_3$
transformer ($\bar d = 0.154$, $n = 25$) and $F_1$ baselines
($\bar d = 0.078$, $n = 25$). The disposition emitted by the verifier
is \texttt{double\_pass\_no\_corpus}.

\begin{table}[h]
\centering
\caption{K2 v2 hypothesis dispositions (Phase 2D, verifier output).
$f$-$p$ is the two-way Type-II ANOVA main-effect $p$-value for the
cell-line-context factor; $p_\text{KW}$ is the Kruskal-Wallis $p$-value
across the three families.}
\label{tab:disp}
\begin{tabular}{lll}
\toprule
Hypothesis & Status & Observed \\
\midrule
$H_1$ capacity & \textbf{PASS} & 4 of 8 datasets at $\rho > 0.5$, $p < 0.05$ \\
$H_{1b}$ data & \textbf{FAIL} & 0 of 8 datasets ($\rho$ undefined per dataset) \\
$H_2$ cell-line & \textbf{PASS} & $f$-$p = 0.0064$, $\eta^2 = 0.120$ \\
$H_{2b}$ corpus & \textbf{FAIL} & per-dataset $\rho > 0$ on every dataset \\
$H_3$ family & \textbf{PASS} & $p_\text{KW} = 1.7\!\times\!10^{-7}$, Cliff's $\delta = 0.80$ \\
\bottomrule
\end{tabular}
\end{table}

\subsection{Ablation panel}
The ablation panel re-evaluates the same hypotheses under counterfactual
analyst-side threshold relaxations. On a frozen Phase 1 snapshot
(\texttt{tests/fixtures/phase1\_results.json}), the locked
$p < 0.05$ threshold passes only one dataset (Replogle~RPE1) for
$H_1$, but the relaxation $\{p < 0.10$ AND $n_\text{pass} \geq 2\}$
admits a second passer (Tahoe, $p = 0.066$) and flips the disposition.
Under the locked thresholds, the flip never occurs. On the Phase 2D
full data, $H_1$ passes under the locked rule (4 of 8 datasets); the
Phase 1 flip case is preserved as evidence of the bias the lock
prevents.

\noindent\emph{K3 PRISM downstream utility.}
Of the 445 ConfPert-calibrated PRISM signatures \cite{corsello2020prism}
that pass BH-FDR, the top-25 ranked by mean coverage deviation include
14 (56\%) with at least one annotated drug-target match in DrugBank or
ChEMBL. The full match table is at
\texttt{results/k3\_tahoe\_predictor\_derived\_top25.json}.

% =============================================================================
% 6. Discussion  (strict-prompt STEP 4: obs + mechanism + comparison + future)
% =============================================================================
\section{Discussion}
\label{sec:discussion}

\paragraph{Observation 1: foundation models are the worst-calibrated
family.}
$F_3$ transformer predictors used in frozen-encoder mode produce
predictive populations with the largest calibration deviation across
all six discrepancies. A plausible mechanism is that the frozen
encoder embeddings, even when followed by a linear head fitted on the
train fold, lose the per-gene variance structure required for matching
the full Perturb-seq distribution; their advantage on supervised
embedding-to-label tasks does not transfer to per-population sampling.
The finding is consistent with \citet{ahlmann2025deep},
\citet{csendes2025train}, and \citet{kedzierska2025zeroshot}, all of
whom report that foundation models lose to a mean baseline on
perturbation prediction under their respective accuracy metrics. The
present paper adds the coverage criterion to that picture but does
\emph{not} settle the full-fine-tune question: we evaluate the
transfer-learning use mode (frozen encoder + ridge head) that
practitioners adopt when downstream labels are scarce, and the result
should be read with that scope in mind.

\paragraph{Observation 2: cell-line context is a real covariate.}
The two-way ANOVA $\eta^2 = 0.120$ on cell-line context exceeds the
pre-registered threshold, indicating that calibration deviation
systematically differs between K562-derived and non-K562 contexts. A
plausible mechanism is that the training corpora for the F2 and F3
predictors are heavily K562-weighted, leading to better calibration on
in-distribution K562 splits than on the more diverse cell lines covered
in Phase 2 (Frangieh melanoma, Schmidt T cells, Datlinger Jurkat).

\paragraph{Observation 3: corpus diversity does not help calibration.}
$H_{2b}$ predicted a negative correlation between log corpus diversity
and mean calibration deviation; the observed sign is positive on every
dataset. One reading is that predictors trained on broader corpora may
overfit to corpus-level structure that does not transfer to held-out
perturbations; another reading is that our corpus-diversity index
(\texttt{predictor.corpus\_diversity\_index} in the YAML) is a poor
proxy. A future experiment would re-fit each predictor on subsamples of
a single source corpus to dissociate diversity from total cell count.

\paragraph{Limitations.}
\emph{Marginal, not conditional coverage.} Distribution-free conditional
coverage is impossible \cite{vovk2012}; ConfPert reports marginal
coverage within head levels.
\emph{Cancer-corpus dominance.} 7 of 8 evaluated datasets are derived
from cancer or transformed cell lines (Lara-Astiaso primary cells is the
exception); the primary-cell shift remains a known confound.
\emph{Heavyweight predictor restriction.} Per pre-reg \S 14.5, the F3
family was evaluated on the 6 heavyweight-allowed datasets only;
McFaline and Lara-Astiaso lack F3 cells. \emph{Frozen-encoder mode.}
The F3 family is evaluated in transfer-learning use, not full fine-tune;
a full \texttt{scgpt.Trainer.train\_perturb} implementation is left to
future work and is expected to narrow but probably not close the gap.
\emph{Power.} $H_1$, $H_{1b}$, $H_{2b}$ are individually underpowered at
their target effect sizes (Appendix~\ref{sec:power}); the headline H3
disposition is well-powered.

\paragraph{Concrete future work.}
(1) Re-evaluate $F_3$ under full perturbation fine-tune
(\texttt{scgpt.Trainer.train\_perturb}) on the same six datasets and
report the calibration gap that closes; an effect-size threshold of
$\Delta \bar d \geq 0.05$ would meaningfully move the $H_3$
disposition. (2) Run the same pre-registered tests on the Replogle
Genome-Wide Perturb-seq (with appropriate compute) to test whether the
H3 PASS replicates on a larger single-dataset sample. (3) Add a CellFM
wrapper once a working PyTorch port appears
(\texttt{biomed-AI/cellfm-800M} returns 404 as of 2026-05-25) and
re-test $H_3$. (4) Replace the corpus-diversity index with a
TF-IDF-style cell-type-coverage measure and re-test $H_{2b}$.

% =============================================================================
% 7. Library and reproducibility  (~1/2 page)
% =============================================================================
\section{Library and reproducibility}
\label{sec:lib}

The benchmark ships as \texttt{confpert} v0.2.0rc2 on PyPI with 126
passing unit tests; a Cell-Eval plug-in covers 5 of 6 discrepancies
absent from the upstream Arc Institute registry; a Codabench
leaderboard at \texttt{paper\_neurips\_dnb\_2026/codabench/} accepts
community submissions. The datasheet (\citealt{gebru2021datasheets};
\texttt{paper\_neurips\_dnb\_2026/datasheet.md}) and model card
(\citealt{mitchell2019model}; \texttt{model\_card.md}) follow the
templates of their respective sources. The Modal volume
\texttt{causeflow-artifacts} holds frozen h5ads and per-job result
JSONs keyed to the locked YAML SHA.

% =============================================================================
% Bibliography
% =============================================================================
\bibliographystyle{plainnat}
\bibliography{../bib/refs}

\appendix

% =============================================================================
% Appendix A: Datasheet for the benchmark (Gebru et al. 2021)
% =============================================================================
\section{Datasheet for ConfPert}
\label{sec:datasheet}

The full datasheet
(\citet{gebru2021datasheets} template) is at
\texttt{paper\_neurips\_dnb\_2026/datasheet.md} and is inlined verbatim
into the camera-ready appendix. It documents motivation, composition
(per-dataset $n_\text{cells}$, perturbation modality, organism,
license), collection process (URLs in \S\ref{sec:datasets}),
preprocessing ($n_\text{hvg} = 512$ HVG selection, \texttt{log1p}
normalisation), uses and tasks, and distribution and maintenance.

% =============================================================================
% Appendix B: Model card (Mitchell et al. 2019)
% =============================================================================
\section{Model card for the confpert library}
\label{sec:modelcard}

The model card is at \texttt{paper\_neurips\_dnb\_2026/model\_card.md}
and is inlined verbatim into the camera-ready appendix. It covers
model details (six conformal heads, one head level), intended use
(calibration coverage benchmarks), factors (predictor family, dataset
context), metrics (six discrepancies), evaluation data (the eight
evaluated datasets), and ethical considerations (cancer-corpus
dominance, primary-cell underrepresentation).

% =============================================================================
% Appendix C: Full K1 results
% =============================================================================
\section{Full K1 results with Wilson 95\% CIs}
\label{sec:k1full}

The per-cell K1 results (651 rows + Wilson 95\% intervals) are stored as
a JSON sidecar at
\texttt{paper\_neurips\_dnb\_2026/phase2d\_report\_phase2.json} under
the key \texttt{bootstrap\_ci\_table}. Each row is keyed by
(predictor, dataset, score, $\alpha$) and reports the held-out
perturbation count, observed coverage $\hat p$, and CI bounds. The
camera-ready will inline a 30-row summary leaderboard; the full table
remains in the JSON to keep the PDF under the 9-page main-body cap.

% =============================================================================
% Appendix D: Verifier output
% =============================================================================
\section{Pre-registration verifier output}
\label{sec:verifier}

CLI reproduction:
\begin{verbatim}
python -m confpert.cli phase2d \
  --prereg paper_neurips_dnb_2026/preregistration_v2.yaml \
  --results baselines/results.json \
  --out paper_neurips_dnb_2026/phase2d_report_phase2.json
\end{verbatim}
Verbatim disposition table is in Table~\ref{tab:disp};
K2 v2 disposition: \texttt{double\_pass\_no\_corpus}.

% =============================================================================
% Appendix E: Power analysis
% =============================================================================
\section{Power analysis (K1 + K2 v2)}
\label{sec:power}

Pre-registered power computations (Phase 2D \texttt{power\_report}):

\begin{tabular}{llll}
\toprule
Test & $n$ & target effect & estimated power \\
\midrule
$H_1$ per-dataset Spearman & 14 & $\rho = 0.5$ & 0.41 (underpowered) \\
$H_{1b}$ cross-dataset Spearman & 10 & $\rho = -0.5$ & 0.25 (underpowered) \\
$H_2$ two-way ANOVA & 120 & $\eta^2 = 0.10$ & 0.81 \\
$H_{2b}$ per-dataset Spearman & 14 & $\rho = -0.5$ & 0.26 (underpowered) \\
$H_3$ Kruskal-Wallis (3 families) & 120 & $\delta = 0.5$ & 0.92 \\
\bottomrule
\end{tabular}

$H_1$, $H_{1b}$, and $H_{2b}$ are individually underpowered at their
target effect sizes; the cross-dataset meta-analysis recovers usable
power for $H_1$ in the headline disposition because the observed
within-dataset effect sizes were larger than the target ($\rho \in
[0.59, 0.71]$ on the passing datasets). $H_2$ and $H_3$ are
independently well-powered.

\end{document}
